% ============================================================================
%  CAPÍTULO III: DESARROLLO DEL TRABAJO
% ============================================================================

\section{Desarrollo del trabajo}

Este capítulo documenta, paso a paso, la configuración de un Controlador de
Dominio de Directorio Activo sobre Debian GNU/Linux 12. Cada bloque describe
desde la preparación del sistema operativo base hasta las pruebas de
autenticación y de unión de clientes al dominio, incluyendo la justificación
técnica de cada una de las decisiones tomadas.

% ============================================================================
\section{Preparación del sistema}

La primera etapa consistió en dejar el servidor listo para asumir el rol de
Controlador de Dominio. Un AD DC es el punto de referencia de identidad de
toda la red, por lo que su nombre y su dirección IP deben ser estables y
conocidos de antemano \citep{msft_activedirectory}.

\subsection{Identidad del host}

A los clientes de la red les llega el DC por \textit{hostname}; si este
cambiara, todos los equipos que ya lo conocen perderían la referencia. Por
ello se fijó el nombre \texttt{dc1} y el nombre completamente cualificado
(FQDN) \texttt{dc1.sudoers.lan}:

\begin{lstlisting}[style=terminal]
$ hostnamectl set-hostname dc1
$ hostnamectl status | head -3
  Static hostname: dc1
  Icon name: computer-laptop
  Operating System: Debian GNU/Linux 12 (bookworm)
\end{lstlisting}

\capturaAPA{fig:hostname}{Identidad de host \texttt{dc1} configurada en el servidor}{captura-hostname.png}{Captura de pantalla mostrando la salida de \texttt{hostnamectl status} con el nombre de host \texttt{dc1}.}

\subsection{Configuración de la IP fija}

Un DC no puede tener una dirección IP variable: Kerberos y el DNS lo
referencian constantemente por \texttt{192.168.0.2}. Además, Samba exige que
el \textit{hostname} resuelva a la IP de la red local y \textbf{nunca} a
\texttt{127.0.0.1} — si resolviera a \textit{loopback}, el DC se ``hablaría
a sí mismo'' y los clientes jamás lo alcanzarían. Por eso se garantizó que
\texttt{dc1} resuelva a \texttt{192.168.0.2/24} en la interfaz cableada
\texttt{eno1} y se registró en \texttt{/etc/hosts}:

\begin{lstlisting}[style=terminal]
$ ip -4 a show eno1
2: eno1: <BROADCAST,MULTICAST,UP,LOWER_UP> ... state UP
    inet 192.168.0.2/24 brd 192.168.0.255 scope global eno1

$ grep -E 'dc1|sudoers' /etc/hosts
192.168.0.2   dc1.sudoers.lan   dc1
\end{lstlisting}

\capturaAPA{fig:ip-fija}{Dirección IP fija \texttt{192.168.0.2/24} en la interfaz \texttt{eno1}}{captura-ip-fija.png}{Captura de pantalla mostrando la salida de \texttt{ip -4 a show eno1} con la IP \texttt{192.168.0.2/24} y la entrada de \texttt{/etc/hosts} que asocia \texttt{dc1.sudoers.lan} a esa dirección.}

\subsection{Instalación de paquetes y servicios del host}

Los paquetes del dominio (Samba 4, cliente Kerberos y utilidades DNS) se
instalaron de forma \textbf{nativa} sobre el host Debian, nunca dentro de un
contenedor: el DC de Samba es un servicio frágil cuya virtualización
contenerizada está descartada por decisión de diseño. La instalación se
realiza con el administrador APT:

\begin{lstlisting}[style=terminal]
$ apt update
$ apt install -y samba smbclient winbind krb5-user \
    libpam-winbind libnss-winbind dnsutils
\end{lstlisting}

Junto con Samba se habilitaron los demás roles del servidor ``todo-en-uno''
que dan sustento físico al dominio: el firewall \texttt{nftables}, el
servidor de hora \texttt{chrony} (imprescindible para Kerberos, que exige un
desfase de reloj menor a 5 minutos) y Docker para los servicios
contenedorizados del resto del proyecto:

\begin{lstlisting}[style=terminal]
$ systemctl is-enabled samba-ad-dc nftables chrony docker
samba-ad-dc.service    enabled
nftables.service       enabled
chrony.service         enabled
docker.service         enabled
\end{lstlisting}

\capturaAPA{fig:servicios-host}{Servicios del host habilitados para el arranque}{captura-servicios-host.png}{Captura de pantalla mostrando la salida de \texttt{systemctl is-enabled} con los servicios \texttt{samba-ad-dc}, \texttt{nftables}, \texttt{chrony} y \texttt{docker} en estado \texttt{enabled}.}

% ============================================================================
\section{Instalación de Samba 4 y promoción del servidor como DC}

Samba 4 integra en un único daemon (\texttt{samba-ad-dc}) todos los
protocolos que forman Active Directory: LDAP (directorio) \citep{rfc4511},
Kerberos (autenticación) \citep{rfc4120}, SMB/CIFS (archivos) y el DNS
propio \citep{sambateam2024}. El servidor se promueve a Controlador de
Dominio mediante el aprovisionamiento del dominio con \texttt{samba-tool}:

\begin{lstlisting}[style=terminal]
$ samba-tool domain provision \
    --use-rfc2307 --realm=SUDOERS.LAN \
    --domain=SUDOERS --server-role=dc \
    --adminpass='AdminPass!2026'
\end{lstlisting}

El aprovisionamiento crea el directorio LDAP, los relatios Kerberos, el
archivo \texttt{/etc/samba/smb.conf} y la maqueta DNS de la zona
\texttt{sudoers.lan}. Tras la promoción, se configura el cliente Kerberos
(\texttt{/etc/krb5.conf}) apuntando al KDC local (\texttt{dc1.sudoers.lan})
y se deja \texttt{samba-ad-dc} como el único demonio del dominio activo,
inhabilitando los heredados \texttt{smbd}, \texttt{nmbd} y \texttt{winbind}
para evitar colisiones de sockets \citep{sambateam2024}.

La promoción correcta se verifica con dos consultas de
\texttt{samba-tool}. La primera muestra el nivel funcional, es decir el
``modo de compatibilidad'' con el que opera el dominio (cuanto más alto,
más funciones activas ofrece):

\begin{lstlisting}[style=terminal]
$ samba-tool domain level show
Domain: SUDOERS
Domain level: 2008 R2
Forest level: 2008 R2
\end{lstlisting}

La segunda pregunta al propio DC \emph{dónde está} y \emph{a qué dominio
sirve}; la línea \texttt{Server role: active directory domain controller}
es la prueba de que la promoción como Controlador de Dominio funcionó:

\begin{lstlisting}[style=terminal]
$ samba-tool domain info 127.0.0.1
Domain:           SUDOERS
Domain SID:       S-1-5-21-...
Forest:           SUDOERS.LAN
Number of users:  13
Server role:      active directory domain controller
\end{lstlisting}

\capturaAPA{fig:domain-info}{Datos del dominio y del Controlador de Dominio}{captura-domain-info.png}{Captura de pantalla mostrando la salida de \texttt{samba-tool domain info 127.0.0.1} con el \textit{realm} \texttt{SUDOERS.LAN}, el nombre NetBIOS \texttt{SUDOERS} y el rol \texttt{active directory domain controller}.}

Finalmente, se comprueba que los puertos del dominio estén escuchando:
\textbf{53} (DNS), \textbf{88} (Kerberos), \textbf{389} (LDAP) y \textbf{445}
(SMB). Es la prueba física de red de que el DC está operativo:

\begin{lstlisting}[style=terminal]
$ ss -ltn | grep -E ':53|:88|:389|:445'
LISTEN  0   128  *:88    *:*   samba  (kerberos)
LISTEN  0   128  *:389   *:*   samba  (ldap)
LISTEN  0   128  *:445   *:*   samba  (smb)
LISTEN  0   128  *:53    *:*   samba  (dns)
\end{lstlisting}

\capturaAPA{fig:puertos-ad}{Puertos del dominio escuchando: 53, 88, 389 y 445}{captura-puertos-ad.png}{Captura de pantalla mostrando la salida de \texttt{ss -ltn} con los puertos 53, 88, 389 y 445 a la escucha.}

% ============================================================================
\section{Configuración del DNS interno}

Un dominio de Active Directory se apoya por completo en el DNS: los clientes
lo consultan para descubrir en qué servidor se encuentra cada servicio de
dominio \citep{rfc1035}. Por ello el DC de Samba asumió la autoridad de la
zona \texttt{sudoers.lan} con el backend \textbf{SAMBA\_INTERNAL} (gestionado
automáticamente por el propio DC, sin competir con el BIND9 que se retiró del
laboratorio).

\subsection{Registro A del Controlador de Dominio}

El registro \texttt{A} es la tabla nombre-$\rightarrow$IP del dominio: toda
consulta empieza pidiendo \texttt{dc1.sudoers.lan}:

\begin{lstlisting}[style=terminal]
$ dig dc1.sudoers.lan
;; ANSWER SECTION:
dc1.sudoers.lan. 900 IN A 192.168.0.2
\end{lstlisting}

\capturaAPA{fig:dig-dc}{Registro A del DC: \texttt{dc1.sudoers.lan} $\rightarrow$ \texttt{192.168.0.2}}{captura-dig-dc.png}{Captura de pantalla mostrando la salida de \texttt{dig dc1.sudoers.lan} con la respuesta \texttt{192.168.0.2}.}

\subsection{Registros SRV de descubrimiento}

Los registros SRV son las ``guías de servicios'' del dominio \citep{rfc2782}:
le dicen a cualquier cliente dónde está el LDAP, el Kerberos o el DC. Sin
ella las estaciones no encuentran el dominio. Su presencia confirma que el
descubrimiento automático funciona sin configuración manual:

\begin{lstlisting}[style=terminal]
$ dig _ldap._tcp.dc._msdcs.sudoers.lan SRV
;; ANSWER SECTION:
_ldap._tcp.dc._msdcs.sudoers.lan. 900 IN SRV 0 100 389 dc1.sudoers.lan.
\end{lstlisting}

\capturaAPA{fig:dig-srv}{Registros SRV de descubrimiento del dominio}{captura-dig-srv.png}{Captura de pantalla mostrando la salida de \texttt{dig \_ldap.\_tcp.dc.\_msdcs.sudoers.lan SRV} con el registro SRV apuntando al puerto 389 de \texttt{dc1.sudoers.lan}.}

\subsection{Zona completa del dominio}

Además de los registros de infraestructura, la zona guarda las direcciones
``lindas'' de los servicios y de los miembros del proyecto
(\texttt{joaquin.sudoers.lan}, \texttt{web.sudoers.lan}, \texttt{print.},
\texttt{mail.}, etc.). Todos resuelven de forma centralizada en el DC:

\begin{lstlisting}[style=terminal]
$ kinit joaquin@SUDOERS.LAN           # ticket para autenticar la consulta
$ samba-tool dns query dc1.sudoers.lan sudoers.lan @ ALL -k yes
Name=, Records=8, Children=0
 DC@SUDOERS.LAN: type=A, value=192.168.0.2
 web@SUDOERS.LAN: type=A, value=192.168.0.2
 joaquin@SUDOERS.LAN: type=A, value=192.168.0.2
...
\end{lstlisting}

\capturaAPA{fig:zona-dns}{Zona DNS completa de \texttt{sudoers.lan} gestionada por el DC}{captura-zona-dns.png}{Captura de pantalla mostrando la salida de \texttt{samba-tool dns query 127.0.0.1 sudoers.lan @ ALL} con los registros A de la zona.}

\subsection{Reenviador hacia Internet}

Para que los clientes resuelvan también nombres externos, el DC funciona con
un \textit{forwarder} hacia el router de la red. De esta manera una única
respuesta DNS del DC cubre tanto los nombres internos (zona
\texttt{sudoers.lan}) como los externos de Internet \citep{iscbind9}.

% ============================================================================
\section{Creación de usuarios y grupos}

Los permisos se otorgan a \textbf{grupos}, no a personas sueltas: de esa
forma ``quién puede qué'' se administra en un solo lugar y los cambios se
propagan a toda la organización. Los grupos y usuarios se crearon con el
script idempotente \texttt{dc/add-users-groups.sh} (puede ejecutarse varias
veces sin duplicar objetos) o manualmente con \texttt{samba-tool}.

\begin{lstlisting}[style=terminal]
$ samba-tool group list | grep -E 'admins|sistemas|oficina|contabilidad'
admins
sistemas
contabilidad
oficina

$ samba-tool user list
joaquin
nicolas
david
grupo2
grupo3
...
\end{lstlisting}

\capturaAPA{fig:grupos}{Grupos de seguridad del dominio}{captura-grupos.png}{Captura de pantalla mostrando la salida de \texttt{samba-tool group list} filtrada por los grupos \texttt{admins}, \texttt{sistemas} y \texttt{oficina}.}

\capturaAPA{fig:usuarios}{Usuarios registrados en el dominio}{captura-usuarios.png}{Captura de pantalla mostrando la salida de \texttt{samba-tool user list} con los usuarios administradores y de prueba.}

Los usuarios son objetos dentro del árbol LDAP, no cuentas sueltas en cada
máquina. Cada objeto tiene un \texttt{sAMAccountName} (el login corto), un
principal de usuario UPN (\texttt{joaquin@SUDOERS.LAN}) y pertenece a
\textit{n} grupos, heredando la suma de sus permisos:

\begin{lstlisting}[style=terminal]
$ samba-tool user show joaquin
dn: CN=joaquin,CN=Users,DC=SUDOERS,DC=LAN
sAMAccountName: joaquin
userPrincipalName: joaquin@SUDOERS.LAN
memberOf: CN=admins,CN=Users,DC=SUDOERS,DC=LAN
          CN=sistemas,CN=Users,DC=SUDOERS,DC=LAN
\end{lstlisting}

\capturaAPA{fig:user-show}{Detalle de un usuario y su pertenencia a grupos}{captura-user-show.png}{Captura de pantalla mostrando la salida de \texttt{samba-tool user show joaquin} con el atributo \texttt{memberOf} enumerando los grupos del usuario.}

La tabla \ref{tab:decisiones} resume las decisiones técnicas más relevantes
de las etapas anteriores y su justificación.

\begin{table}[H]
\centering
\caption{Decisiones técnicas y justificación}
\label{tab:decisiones}
\contenidoConFuente{%
  \begin{tablaAPA}
    \begin{tabular}{@{}p{5cm}p{8.5cm}@{}}
    \toprule
    \textbf{Decisión} & \textbf{Justificación} \\
    \midrule
    DC nativo (fuera de Docker) & Samba 4 como DC es frágil en contenedores; se despliega directamente sobre Debian \citep{sambateam2024} \\
    IP fija \texttt{192.168.0.2} / hostname \texttt{dc1} & Kerberos y DNS referencian constantemente al DC; un cambio de IP o nombre rompe el dominio \\
    \texttt{samba-tool domain provision} & Método oficial de promoción a DC; genera configuración, directorio LDAP y maqueta DNS \\
    DNS \textbf{SAMBA\_INTERNAL} autoritativo & El AD exige que el DC sea quien sirva la zona \texttt{sudoers.lan} y los registros SRV \\
    Permisos por grupos AD & Administración centralizada de ``quién puede qué''; herencia automática de permisos \\
    \texttt{samba-ad-dc} como único demonio & Evita colisiones de puertos con \texttt{smbd}, \texttt{nmbd} y \texttt{winbind} heredados \\
    \bottomrule
    \end{tabular}
  \end{tablaAPA}%
}{Elaboración propia.}
\end{table}

% ============================================================================
\section{Pruebas de autenticación}

El protocolo Kerberos resuelve el problema de ``demostrar quién soy sin
enviar la contraseña por la red'' \citep{rfc4120}. La respuesta son
\textbf{tickets}: el KDC (el propio DC) emite un ticket de concesión (TGT) al
autenticar al usuario, y ese ticket viaja firmado por el DC, de modo que
nadie puede falsificar una identidad.

La prueba en el servidor se realiza con las herramientas del paquete
\texttt{krb5-user}:

\begin{lstlisting}[style=terminal]
$ kinit juan@SUDOERS.LAN
Password for juan@SUDOERS.LAN:
$ klist
Ticket cache: FILE:/tmp/krb5cc_1000
Default principal: juan@SUDOERS.LAN
Valid starting     Expires            Service principal
09/20/26 10:00:00  09/20/26 20:00:00  krbtgt/SUDOERS.LAN@SUDOERS.LAN
$ kdestroy
\end{lstlisting}

\capturaAPA{fig:kinit-klist}{Solicitud y verificación de ticket Kerberos (TGT)}{captura-kinit-klist.png}{Captura de pantalla mostrando la salida de \texttt{kinit} y \texttt{klist} sobre \texttt{dc1}, con el ticket emitido por el KDC \texttt{SUDOERS.LAN} y su vigencia.}

\subsection{Autenticación desde un cliente real}

La misma cuenta y el mismo ``carnet'' funcionan desde cualquier equipo de la
red, porque el emisor de tickets siempre es el DC. Esto es el
*single sign-on*: el cliente se identifica una vez y accede a los servicios
del dominio \citep{msft_activedirectory}.

\begin{lstlisting}[style=terminal]
$ sudo kinit juan@SUDOERS.LAN && klist   # desde el cliente unido
Default principal: juan@SUDOERS.LAN
Valid starting     Expires            Service principal
09/20/26 10:05:00  09/20/26 20:05:00  krbtgt/SUDOERS.LAN@SUDOERS.LAN
\end{lstlisting}

\capturaAPA{fig:kinit-cliente}{Autenticación Kerberos desde un equipo cliente}{captura-kinit-cliente.png}{Captura de pantalla ejecutando \texttt{kinit} y \texttt{klist} en el equipo cliente, mostrando el ticket obtenido contra el DC.}

% ============================================================================
\section{Unión de clientes al dominio}

Unirse al dominio es más que configurar el DNS \citep{msft_activedirectory}:
\texttt{realm join} crea una \textbf{computer account} en el AD, convirtiendo
al equipo en un objeto del dominio con su propio secreto de máquina. Desde
entonces el cliente usa el AD como fuente de identidad (LDAP para leer,
Kerberos para autenticar) sin copiar usuarios locales.

\subsection{Unión del cliente Debian}

El cliente se unió al dominio ejecutando \texttt{realm join --user=administrator
SUDOERS.LAN} desde el equipo, autenticándose con la cuenta de administración
del dominio. La operación fue exitosa y quedó registrada la \emph{computer
account} correspondiente, como se verifica en el servidor más adelante.

\subsection{Verificación desde el cliente}

Los usuarios del AD aparecen como si fueran locales gracias a SSSD:
\texttt{getent} los consulta en vivo contra el DC, sin duplicar cuentas en el
equipo.

\begin{lstlisting}[style=terminal]
$ getent passwd joaquin david nicolas
joaquin:*:10001:10001:...:/home/joaquin:/bin/bash
david:*:10002:10002:...:/home/david:/bin/bash
nicolas:*:10003:10003:...:/home/nicolas:/bin/bash

$ getent group 'admins'
admins:*:20001:joaquin,david,nicolas
\end{lstlisting}

\capturaAPA{fig:getent}{Usuarios del dominio visibles desde el cliente con \texttt{getent}}{captura-getent.png}{Captura de pantalla mostrando la salida de \texttt{getent passwd} y \texttt{getent group} con los usuarios del AD resueltos en el cliente.}

\subsection{Verificación desde el servidor}

Del lado del server se comprueba que la notebook haya quedado registrada como
equipo del dominio:

\begin{lstlisting}[style=terminal]
$ samba-tool computer list
DESKTOP-NBK1$
...
\end{lstlisting}

\capturaAPA{fig:computer-list}{Estaciones registradas en el dominio}{captura-computer-list.png}{Captura de pantalla mostrando la salida de \texttt{samba-tool computer list} con el equipo del cliente unido al AD.}

\subsection{Acceso a recursos de archivo}

El DC sirve shares de forma nativa (sin contenedor) con permisos definidos
por grupos AD en el servidor: \texttt{departamentos}, \texttt{homes} (cada
usuario ve su propio directorio) y \texttt{respaldo}. El acceso se verifica
desde el servidor con \texttt{smbclient} y desde el cliente mediante un
montaje CIFS o el gestor de archivos:

\begin{lstlisting}[style=terminal]
$ smbclient -L dc1 --user=joaquin          # en el server
	Sharename       Type      Comment
	---------       ----      -------
	departamentos   Disk      Recursos por departamentos
	homes           Disk      Directorios personales
	respaldo        Disk
\end{lstlisting}

\capturaAPA{fig:smbclient}{Listado de los recursos compartidos (SMB) del DC}{captura-smbclient.png}{Captura de pantalla mostrando la salida de \texttt{smbclient -L dc1 --user=joaquin} con los shares \texttt{departamentos}, \texttt{homes} y \texttt{respaldo}.}

% ============================================================================
\section{DHCP y DNS en las estaciones de la red}

Para que una estación recién conectada quede ``lista para el dominio'' sin
intervención manual, la red cuenta con el servidor DHCP Kea (en contenedor,
con \texttt{network\_mode: host} porque DHCP usa broadcast y no atraviesa un
bridge de Docker). Kea entrega junto con la IP las \textbf{opciones} de red:
DNS \texttt{192.168.0.2}, dominio \texttt{sudoers.lan}, gateway
\texttt{192.168.0.1} y NTP \texttt{192.168.0.2}. El DHCP del router está
desactivado para evitar el doble servidor DHCP.

\subsection{Verificación del servicio DHCP}

\begin{lstlisting}[style=terminal]
$ sudo ss -ulnp | grep ':67 '
UNCONN  0  0   192.168.0.2:67  *:*  users:(("kea-dhcp4",pid=...))

$ sudo cat /var/lib/kea/kea-leases4.csv   # concesiones en vivo
address,hostname,state,valid_lifetime,...
192.168.0.100,notebook1,assigned,...
\end{lstlisting}

\capturaAPA{fig:kea-ss}{Servidor DHCP Kea escuchando en \texttt{192.168.0.2:67}}{captura-kea-ss.png}{Captura de pantalla mostrando la salida de \texttt{ss -ulnp} con Kea a la escucha en el puerto 67/udp y las concesiones activas.}

\subsection{Prueba de obtención de IP desde un cliente}

Un cliente se desconecta y vuelve a tomar IP, quedando con la configuración
completa entregada por Kea:

\begin{lstlisting}[style=terminal]
$ sudo dhclient -v eno2
DHCPDISCOVER on eno2 to 255.255.255.255 port 67
DHCPOFFER from 192.168.0.2
DHCPACK from 192.168.0.2 (xid=0x...)
$ ip -4 a
inet 192.168.0.101/24 ...   # DNS = 192.168.0.2, dominio = sudoers.lan
$ dig dc1.sudoers.lan      # resuelve con el DNS entregado por DHCP
\end{lstlisting}

\capturaAPA{fig:dhclient}{Cliente tomando IP del pool de Kea y resolviendo el DC}{captura-dhclient.png}{Captura de pantalla mostrando \texttt{dhclient -v} obteniendo una IP del pool dinámico y la posterior consulta \texttt{dig dc1.sudoers.lan} resuelta con el DNS entregado por DHCP.}